% Generated draft: requires mathematical and source review.
\documentclass[11pt]{article}
\usepackage[margin=25mm]{geometry}
\usepackage{fontspec}
\setmainfont{Latin Modern Roman}
\usepackage{amsmath,amssymb,mathtools,unicode-math}
\setmathfont{Latin Modern Math}
\usepackage{enumitem,needspace,xurl,hyperref}
\setlength{\parindent}{0pt}
\setlength{\parskip}{5pt}
\setlength{\emergencystretch}{3em}
\newcommand{\N}{\mathbb{N}}
\newcommand{\Z}{\mathbb{Z}}
\newcommand{\Q}{\mathbb{Q}}
\newcommand{\R}{\mathbb{R}}
\newcommand{\C}{\mathbb{C}}
\newcommand{\F}{\mathbb{F}}
\newcommand{\A}{\mathbb{A}}
\newcommand{\PP}{\mathbb{P}}
\newcommand{\E}{\mathbb{E}}
\newcommand{\eps}{\varepsilon}
\newcommand{\dd}{\,\mathrm{d}}
\begin{document}
\section*{UnsolvedMath ID: 30002031}
\textbf{Source problem code:} catalogue entry
\subsection{Definitions and mathematical statement}
Let \(r\ge 2\) and \(n\ge r\) be integers, let \([n]=\{1,\ldots,n\}\), and write \(\binom{X}{k}\) for the family of all \(k\)-element subsets of a finite set \(X\). For disjoint sets \(P,Q\) with \(|P|=r-2\) and \(|Q|=4\), define the \emph{\(r\)-daisy} by
\[
D_r(P,Q)=\{P\cup B:B\in\tbinom{Q}{2}\}.
\]
This is an \(r\)-uniform hypergraph with six edges on \(r+2\) vertices; \(D_r\) denotes its isomorphism type. In particular, \(P=\varnothing\) is allowed when \(r=2\).

A family \(\mathcal A\subseteq\binom{[n]}{r}\) is \(D_r\)-free if it contains no copy of these six edges. Explicitly, for every \(P\in\binom{[n]}{r-2}\) and every \(Q\in\binom{[n]\setminus P}{4}\), there is a \(B\in\binom{Q}{2}\) such that \(P\cup B\notin\mathcal A\). Containment here is ordinary hypergraph containment, with no induced-subhypergraph requirement.

Define
\[
\operatorname{ex}(n,D_r)=\max\bigl\{|\mathcal A|:\mathcal A\subseteq\tbinom{[n]}{r},\ \mathcal A\text{ is }D_r\text{-free}\bigr\}.
\]
The problem asks to determine how large such a family can be, equivalently to determine \(\operatorname{ex}(n,D_r)\) as a function of \(n\) and \(r\). When \(r\le n<r+2\), the forbidden configuration cannot fit on the ground set, so the condition is vacuous and the maximum is \(\binom nr\).

For the density questions mentioned in the supplied title and literature assessment, the associated Turán density is
\[
\pi(D_r)=\lim_{n\to\infty}\frac{\operatorname{ex}(n,D_r)}{\binom nr},
\]
where \(r\) is fixed in this limit. The source question does not specify whether an exact finite-\(n\) formula, an asymptotic determination, or bounds are the intended target.

The supplied dated curation labels the problem partially solved and attributes to Ellis, Ivan, and Leader a disproof of the conjecture \(\lim_{r\to\infty}\pi(D_r)=0\). It also reports that the extremal-function and exact-density questions, including the value of \(\pi(D_3)\), remain open. These are supplied literature assessments, not independently verified assertions in this draft.
\subsection{Short English statement}
For each integer \(r\ge2\), an \(r\)-daisy consists of all six ways to add two elements of a fixed four-element set to a disjoint fixed \((r-2)\)-element set. Determine the maximum number of \(r\)-element subsets of an \(n\)-element set that can be chosen without containing all six members of any such daisy. The source also discusses the associated asymptotic densities; its partially solved status is supplied dated curation.
\subsection{Sources}
\begin{itemize}
\item Local UnsolvedMath registry, supplied record with permanent ID \(30002031\), title \emph{Turán Densities of Hypergraph Daisies}, in the Oberwolfach Reports Open Problems collection. Exact source problem code: \texttt{catalogue entry}. No individual proposer is identified in the record.
\item Supplied report citation: \emph{Mini-Workshop: Hypergraph Turán Problem (2013)}, Oberwolfach Reports, pp.\ 1149--1182, DOI \texttt{10.4171/owr/2012/19}. \url{https://doi.org/10.4171/owr/2012/19}. The registry separately records proposed year \(2012\); the report-title year and the year appearing in the DOI are retained as supplied without resolving their discrepancy.
\item Supplied literature reference labeled Ellis--Ivan--Leader: \url{https://arxiv.org/abs/2401.16289}. The record attributes the disproof of the vanishing-daisy-density conjecture to Ellis, Ivan, and Leader.
\item Supplied literature reference labeled Daisy density results: \url{https://api.repository.cam.ac.uk/server/api/core/bitstreams/11513fbb-eac3-4eaf-bed3-46228c0c6758/content}.
\item Supplied literature reference labeled Oberwolfach report: \url{https://ems.press/content/serial-article-files/46390?nt=1}.
\item Literature assessment checked at \texttt{2026-08-21}, exactly as recorded. The partially solved label, attributed disproof, and statements about remaining open questions are supplied dated curation derived from the OWR collection, not independently verified here. The supplied background explicitly requests independent verification of time-sensitive status and cited claims.
\end{itemize}
\end{document}
